\subsection{Simulations for Orchestration Performance of Workload Predictors}
\label{sec:ch3_simulations}

% https://ieeexplore.ieee.org/abstract/document/9796778

To evaluate the impact of workload prediction on container orchestration performance, we selected \ac{COSCO}~\cite{tuli2021cosco} for its validated accuracy and strong correlation between simulated and actual fog environment metrics. The framework implements a discrete-time control system that tracks \ac{CPU}, \ac{RAM}, disk, and network utilization across simulated hosts at fixed intervals. For our evaluation, we configured a heterogeneous fog computing environment with both edge and cloud resources, and contributed a prediction interface to enable systematic assessment of workload prediction impact without physical deployment costs.

\subsubsection{Experimental Configuration}

We evaluate our scheduling approach using the \ac{COSCO} framework~\cite{tuli2021cosco}, which simulates realistic edge-cloud infrastructures. Their deep learning \texttt{GOBIScheduler} orchestrates containers across the distributed infrastructure $\mathcal{H}$, consisting of $|\mathcal{H}| = 50$ heterogeneous computational nodes with a total power capacity of $1.0\,\text{kW}$ and inter-node network bandwidth of $10\,\text{Gb}/\text{s}$. The simulation executes $2000$ discrete time steps at $300\,\text{ms}$ intervals, with $|\mathcal{C}| = 50$ containers and $4$ new container introductions per interval.

 We use the \texttt{AzureFog} configuration with 60\% edge nodes and 40\% cloud nodes. The nodes consist of three \ac{VM} types: B2s \acp{VM} (40\% of nodes) with 4,029 \ac{MIPS}, 4GB \ac{RAM}, and 32GB storage. B4ms \acp{VM} (40\%) have 4,029 \ac{MIPS}, 16GB \ac{RAM}, and 32GB storage. B8ms \acp{VM} (20\%) have 16,111 \ac{MIPS}, 32GB \ac{RAM}, and 32GB storage.

Edge nodes have a network latency of \SI{3}{\milli\second}, while cloud nodes operate at \SI{76}{\milli\second}. All nodes are equipped with 5Gbps network bandwidth for both uplink and downlink. Each node $h_i \in \mathcal{H}$ has varying capacities for \ac{CPU}, memory, disk, and network resources, as defined in Section~\ref{sec:ch3_system_model}. Each container $c \in \mathcal{C}$ comprises tasks $\mathcal{J}_c$, where each task $j \in \mathcal{J}_c$ arrives at time $a_j$ and requires resources $\mathbf{r}_j$.

\subsubsection{Dual-Criteria Orchestration Objective Function}

The \texttt{GOBIScheduler} minimizes a normalized dual-objective function over each interval $[t, t+1]$ that balances energy efficiency and response time \cite{tuli2021cosco}, two of the most important metrics for edge-cloud computing \cite{mcchesney2019defog}:

\begin{equation}
\label{eq:objective_function}
\mathcal{O}(\mathbf{P}_t) = \alpha \, \mathcal{E}_t^{\text{norm}} + (1-\alpha) \, \mathcal{R}_t^{\text{norm}},
\end{equation}

\noindent where $\alpha \in [0,1]$ controls the trade-off between energy consumption and response time, and $\mathbf{P}_t$ represents the system state at time $t$. The normalized energy consumption $\mathcal{E}_t^{\text{norm}}$ measures the total power used across all nodes relative to their maximum capacity:

\begin{equation}
\label{eq:energy_consumption}
\mathcal{E}_t^{\text{norm}} = \frac{\mathcal{E}_t}{\mathcal{E}_t^{\text{max}}} = \frac{\sum_{h_i \in \mathcal{H}} \int_{t}^{t+1} \rho_{h_i}(s) \, ds}{|\mathcal{C}_t| \sum_{h_i \in \mathcal{H}} \rho_{h_i}^{\text{max}} \Delta t},
\end{equation}

\noindent where $\rho_{h_i}(s)$ is the instantaneous power consumption of node $h_i$ at time $s$, $\rho_{h_i}^{\text{max}}$ is its maximum power capacity, $\mathcal{C}_t$ is the set of active containers, and $\Delta t = 1$ is the interval duration. The normalized response time $\mathcal{R}_t^{\text{norm}}$ quantifies task completion efficiency across the container set $\mathcal{C}$:

\begin{equation}
\label{eq:response_time}
\mathcal{R}_t^{\text{norm}} = \frac{\mathcal{R}_t}{\mathcal{R}_t^{\text{max}}} = \frac{\sum_{j \in \mathcal{J}_t^{\text{complete}}} \Delta t_j}
{|\mathcal{J}_t^{\text{complete}}| \cdot \max_{s \leq t} \Delta t_s},
\end{equation}

\noindent where $\Delta t_j = t_j^{\text{end}} - t_j^{\text{start}}$ represents the completion time for task $j$.

\subsubsection{Integration of MLFlow-Based Orchestration}

\begin{figure}[!htbp]
\centering
\includegraphics[width=0.568\textwidth]{sections/chapter3/img/simulation_flowchart.pdf}
\caption{System architecture of the extended \ac{COSCO} simulator. The \ac{MLFlow} Registry provides pre-trained models to the Predictor component, which works alongside the workload generator to inform scheduling decisions. The Environment maintains the system state while metrics are continuously collected.}
\label{fig:simulation_flowchart}
\end{figure}

To evaluate our proposed \ac{AERO} model (Section~\ref{sec:ch3_proposed_solution}) within a simulated environment, we extended the \ac{COSCO} framework~\cite{tuli2021cosco} with a prediction interface integrated with \ac{MLFlow}~\cite{MLFlow}, as illustrated in Fig.~\ref{fig:simulation_flowchart}. The novel simulator operates in two modes: the traditional \textit{Reactive Mode}, where the \ac{BitBrains} workload generator directly informs the scheduler of current resource demands, and our novel \ac{MLFlow} Mode, where pre-trained models from the \ac{MLFlow} Registry are used through the prediction interface to forecast future resource demands.


\subsubsection{BitBrains Workload Dataset Analysis and Model Training}
\label{sec:ch3_bitbrains}

\begin{figure*}[!htbp]  
\centering
\includegraphics[width=0.9\textwidth]{sections/chapter3/img/feature_grid_comparison.pdf}
\caption{Container resource utilization predictions and normalized \ac{RMSE} comparison across different metrics, with bars ordered by performance.}
\label{fig:feature_grid_comparison}
\end{figure*}

We use the \ac{BitBrains} Random trace~\cite{BitbrainsDS} for our simulation evaluations. The dataset consists of traces from 500 virtual machines connected to \ac{SAN} and \ac{NAS} devices. Each of the 500 files provides $T = 8,631$ samples collected at \SI{300}{\milli\second} intervals. Each workload vector $\mathbf{W}_t \in \mathbb{R}^{f}$ consists of $f = 7$ features, as described in Table~\ref{tab:predicted_features}, aligning with the workload representation in Section~\ref{sec:ch3_system_model}.

\begin{table}
\centering
\renewcommand{\arraystretch}{1.15}
\begin{tabular}{p{0.35\textwidth}p{0.55\textwidth}}
\hline
\textbf{Feature} & \textbf{Description} \\
\hline
$\mathrm{CPU\;usage\;[MHz]}$ & Raw \ac{CPU} computational power \\
$\mathrm{CPU\;usage\;[\%]}$ & \ac{CPU} utilization percentage \\
$\mathrm{Memory\;usage\;[KB]}$ & \ac{RAM} consumption \\
$\mathrm{Disk\;read\;[KB/s]}$ & Disk read throughput \\
$\mathrm{Disk\;write\;[KB/s]}$ & Disk write throughput \\
$\mathrm{Network\;recv.\;[KB/s]}$ & Incoming network throughput \\
$\mathrm{Network\;trans.\;[KB/s]}$ & Outgoing network throughput \\
\hline
\end{tabular}
\caption{Key resource usage features being predicted by the models}
\label{tab:predicted_features}
\end{table}

Following the methodology in Section~\ref{sec:ch3_evaluation}, we trained \ac{AERO} and compared its performance with Pathformer and \ac{SparseTSF} on the \ac{BitBrains} dataset. Fig.~\ref{fig:feature_grid_comparison} shows the \ac{RMSE} comparisons across various resource metrics. The bar chart displays normalized \ac{RMSE} values for cross-metric comparison, while absolute \ac{RMSE} values in original units are reported in the text below. Consistent with previous findings, \ac{AERO} performs well across most resource dimensions.

\ac{AERO} outperforms the baselines in five out of seven metrics, with significant improvements in \ac{CPU} utilization prediction (\ac{RMSE} of 2.30\,\ac{MHz} vs.\ 5.96\,\ac{MHz} for Pathformer) and network receive throughput (0.12\,KB/s vs.\ 0.42\,KB/s). This advantage stems from \ac{AERO}'s ability to dynamically adapt to workload periodicities (Section~\ref{sec:ch3_dynamic_periodicity}), whereas \ac{SparseTSF}'s fixed periodicity assumption limits its adaptability.

While \ac{SparseTSF} slightly outperforms \ac{AERO} in memory usage prediction (\ac{RMSE} of 25.54\,MB compared to 26.09\,MB for \ac{AERO}), \ac{AERO} shows consistent performance across multiple metrics. Combined with its lightweight architecture (Section~\ref{sec:ch3_model_complexity}), this shows its suitability for real-time cloud orchestration in highly dynamic environments.

\subsubsection{Impact of Workload Predictions on Orchestration Performance}

\begin{figure}[!htbp]
    \centering
    \begin{subfigure}[b]{0.32\textwidth}
        \centering
        \includegraphics[width=\linewidth]{sections/chapter3/img/2x4_metrics_grid/energy_consumption.pdf}
        \caption{Energy Consumption $\mathcal{E}_t$}
        \label{subfig:energy}
    \end{subfigure}
    \hfill
    \begin{subfigure}[b]{0.32\textwidth}
        \centering
        \includegraphics[width=\linewidth]{sections/chapter3/img/2x4_metrics_grid/response_time.pdf}
        \caption{Response Time $\mathcal{R}_t$}
        \label{subfig:response}
    \end{subfigure}
    \hfill
    \begin{subfigure}[b]{0.32\textwidth}
        \centering
        \includegraphics[width=\linewidth]{sections/chapter3/img/2x4_metrics_grid/cpu_utilization.pdf}
        \caption{\ac{CPU} Utilization}
        \label{subfig:cpu}
    \end{subfigure}
    \\[1ex]
    \begin{subfigure}[b]{0.32\textwidth}
        \centering
        \includegraphics[width=\linewidth]{sections/chapter3/img/2x4_metrics_grid/ram_utilization.pdf}
        \caption{\ac{RAM} Utilization}
        \label{subfig:ram}
    \end{subfigure}
    \hfill
    \begin{subfigure}[b]{0.32\textwidth}
        \centering
        \includegraphics[width=\linewidth]{sections/chapter3/img/2x4_metrics_grid/sla_violations.pdf}
        \caption{\ac{SLA} Violations}
        \label{subfig:sla}
    \end{subfigure}
    \hfill
    \begin{subfigure}[b]{0.32\textwidth}
        \centering
        \includegraphics[width=\linewidth]{sections/chapter3/img/2x4_metrics_grid/migration_counts.pdf}
        \caption{Task Migrations}
        \label{subfig:migration_count}
    \end{subfigure}
    \caption{Comparison of resource utilization, energy consumption, and service quality metrics among \ac{AERO}, Pathformer, \ac{SparseTSF}, and Reactive scheduling approaches.}
    \label{fig:metrics_comparison}
\end{figure}

To assess how workload predictions influence the orchestration objectives defined in Section~\ref{sec:ch3_simulations}, we compared our \ac{AERO} model with other predictive schedulers (Pathformer, \ac{SparseTSF}) and a Reactive baseline that does not use predictions. As shown in Fig.~\ref{fig:metrics_comparison}, we evaluated key performance metrics including energy consumption $\mathcal{E}_t$, response time $\mathcal{R}_t$, resource utilization, and \ac{SLA} violations, based on the dual-objective function in Equation~\eqref{eq:objective_function}.

\paragraph*{Energy Consumption and \ac{CPU} Utilization}

Fig.~\ref{fig:metrics_comparison} compares the energy consumption $\mathcal{E}_t$ and \ac{CPU} utilization across different models. The Reactive approach consumes more energy, with a mean $\mathcal{E}_t$ of $1293$\,J, compared to $1123$\,J for both \ac{AERO} and Pathformer. This higher energy usage is directly related to increased \ac{CPU} utilization, as shown in Fig.~\ref{subfig:cpu}.

Without workload predictions, the Reactive model cannot anticipate resource demands, leading to delayed scaling. As workloads spike unexpectedly, \ac{CPU} utilization surges, causing higher energy consumption according to Equation~\eqref{eq:energy_consumption}. Predictive models like \ac{AERO} forecast future demands, allowing the scheduler to adjust resources proactively. This foresight results in smoother \ac{CPU} usage and reduced energy consumption, optimizing the objective $\mathcal{O}(\mathbf{P}_t)$.

\paragraph*{Response Time and \ac{SLA} Violations}

Timely resource provisioning also impacts response time $\mathcal{R}_t$ (Equation~\eqref{eq:response_time}) and \ac{SLA} violations. The Reactive model experiences higher response times, with a mean $\mathcal{R}_t$ of $10.02$\,s, compared to $3.29$\,s for \ac{AERO} and $2.45$\,s for Pathformer (Fig.~\ref{subfig:response}). Late adjustments to workload changes mean tasks wait longer, increasing $\mathcal{R}_t$ and negatively affecting $\mathcal{O}(\mathbf{P}_t)$.

Additionally, the Reactive approach shows a higher \ac{SLA} violation rate (Fig.~\ref{subfig:sla}) at $22.21\%$, whereas \ac{AERO} and Pathformer maintain low rates of $0.21\%$ and $0.10\%$, respectively. By predicting workload spikes, \ac{AERO} ensures resources are available when needed, minimizing response times and \ac{SLA} violations, thus improving overall orchestration performance.

\paragraph*{Task Migrations and Resource Utilization}

As shown in Fig.~\ref{subfig:migration_count}, the Reactive model performs the fewest task migrations (2,026), while the predictive models exhibit more migrations: \ac{SparseTSF} (5,309), \ac{AERO} (7,747), and Pathformer (8,480). This increase is due to proactive workload predictions enabling the scheduler to reallocate tasks in anticipation of resource demand spikes.

Such migrations enable the system to adapt to dynamic workload conditions, ensuring compliance with \acp{SLA} and improving overall performance. By proactively redistributing tasks, the scheduler prevents resource overloads and balances utilization across the network. Moreover, in scenarios where workloads are low across different servers, migrations facilitate server consolidation by moving tasks to shared machines, resulting in significant energy savings. As indicated in Fig.~\ref{subfig:cpu}, predictive models maintain steadier \ac{CPU} utilization compared to the Reactive model, which suffers from resource hotspots due to its inability to anticipate workload spikes.

% The BitBrains workload dataset~\cite{BitbrainsDS} is characterized by spiky and volatile behavior in CPU, memory, disk, and network transmissions (Fig.~\ref{fig:feature_grid_comparison}). Such high variability causes resource demands to change rapidly and unpredictably. Predictive models like AERO anticipate these sudden fluctuations, prompting proactive task migrations to prevent overloads and resource contention.

\paragraph*{Overall Implications}

Incorporating workload predictions improves the dual-objective function (Equation~\eqref{eq:objective_function}). Predictive models like \ac{AERO} enable proactive resource management, reducing both energy consumption and response time. In dynamic fog computing environments, predictive orchestration optimizes performance metrics.

While accurate predictions are beneficial, the adaptability and efficiency of the orchestration model are more critical for improving $\mathcal{O}(\mathbf{P}_t)$. \ac{AERO} balances prediction accuracy with computational efficiency, making it suitable for real-time deployments without excessive overhead.
